% ============================================================
%  Chapter 4 -- Limits of Functions, 4.1-4.4
% ============================================================
\rchapter{4}{Continuity}

\begin{roadmap}[Why this chapter exists]
Chapters 2 and 3 built machinery --- topology and limits --- and this chapter
is where it gets spent. Continuity is defined three equivalent ways
($\varepsilon$--$\delta$ at 4.5, sequences via 4.2, inverse images at 4.8),
and then the two great theorems of Chapter 2 are converted into the two
facts every calculus student uses without proof: a continuous function on a
compact set is bounded, attains its extrema, and is uniformly continuous
(4.14--4.19); a continuous function on an interval takes every intermediate
value (4.22--4.23).

\medskip
The chapter's centre of gravity is \textbf{Theorem 4.8}: continuous means
\emph{inverse images of open sets are open}. It converts continuity from a
pointwise, quantifier-laden condition into a single structural statement, and
it is the form in which every later argument uses the hypothesis.
\end{roadmap}

\begin{roadmap}[How it is organized, and why in that order]
\textbf{Rudin's own account}, from the opening paragraph: the theorems ``would
not become any easier if we restricted ourselves to real functions,'' and it
``actually simplifies and clarifies the picture to discard unnecessary
hypotheses and to state and prove theorems in an appropriately general
context.'' So everything is done for metric spaces, exactly as in Chapters 2
and 3, and specialised only when the real line genuinely matters
(monotonicity, 4.28--4.30, needs an order; the intermediate value theorem
needs intervals).

\medskip
\centering
\begin{tikzpicture}[scale=0.9, transform shape, node distance=4mm]
\node[rmbox, text width=19mm] (a) {\textbf{4.1--4.12}\\limits;\\continuity\\three ways};
\node[rmbox, text width=19mm, right=of a] (b) {\textbf{4.13--4.21}\\compactness\\pays out};
\node[rmbox, text width=19mm, right=of b] (c) {\textbf{4.22--4.24}\\connectedness\\pays out};
\node[rmbox, text width=19mm, right=of c] (d) {\textbf{4.25--4.34}\\discontinuities;\\monotonicity};
\draw[rmarrow] (a) -- (b); \draw[rmarrow] (b) -- (c); \draw[rmarrow] (c) -- (d);
\end{tikzpicture}

\medskip
\raggedright
Three things to know in advance.

\smallskip
\textbf{The proofs are translations, not inventions.} Almost every proof here
either translates a sequence theorem through the bridge 4.2 (so Chapter 3 does
the work), or pulls a topological property of the target space back through
$f^{-1}$ (so Chapter 2 does the work). The strategy index of Chapter 1 gains
an eighth pattern: \emph{transport structure along $f$ or $f^{-1}$}.

\smallskip
\textbf{The counterexamples are half the content.} 4.20--4.21 show exactly
what dies without compactness, and 4.24, 4.27 what the intermediate value
property does \emph{not} imply. Rudin places a counterexample after nearly
every theorem in this chapter; read them as part of the theorem.

\smallskip
\textbf{Uniform continuity (4.18--4.19) is the new concept.} Everything else
repackages Chapters 2--3; uniformity --- one $\delta$ serving every point at
once --- is genuinely new, is the hardest distinction in the chapter, and is
the hypothesis integration will need in Chapter 6 (Theorem 6.8).
\end{roadmap}

\noindent
The function concept and some of the related terminology were introduced in
Definitions 2.1 and 2.2. Although we shall (in later chapters) be mainly
interested in real and complex functions (i.e., in functions whose values are
real or complex numbers) we shall also discuss vector-valued functions (i.e.,
functions with values in $\R^k$) and functions with values in an arbitrary
metric space. The theorems we shall discuss in this general setting would not
become any easier if we restricted ourselves to real functions, for instance,
and it actually simplifies and clarifies the picture to discard unnecessary
hypotheses and to state and prove theorems in an appropriately general
context.

The domains of definition of our functions will also be metric spaces,
suitably specialized in various instances.

\rsection{Limits of Functions}

\ritem{4.1}{Definition}
Let $X$ and $Y$ be metric spaces; suppose $E \subset X$, $f$ maps $E$ into
$Y$, and $p$ is a limit point of $E$. We write $f(x) \to q$ as $x \to p$, or
\begin{equation}
  \lim_{x\to p} f(x) = q \tag{4.1}
\end{equation}
if there is a point $q \in Y$ with the following property: For every
$\varepsilon > 0$ there exists a $\delta > 0$ such that
\begin{equation}
  d_Y(f(x),q) < \varepsilon \tag{4.2}
\end{equation}
for all points $x \in E$ for which
\begin{equation}
  0 < d_X(x,p) < \delta. \tag{4.3}
\end{equation}
The symbols $d_X$ and $d_Y$ refer to the distances in $X$ and $Y$,
respectively.\kw{Two structural points hide in (4.3). The strict inequality
$0 < d_X(x,p)$ \emph{deletes the point $p$ itself}: the limit ignores what
happens at $p$, and even whether $f$ is defined there. And $p$ must be a
\emph{limit point} of $E$ --- otherwise no $x$ satisfies (4.3) for small
$\delta$ and every $q$ would qualify vacuously. The hypothesis is exactly
what makes the limit unique.}

If $X$ and/or $Y$ are replaced by the real line, the complex plane, or by some
euclidean space $\R^k$, the distances $d_X$, $d_Y$ are of course replaced by
absolute values, or by norms of differences (see Example 2.16).

It should be noted that $p \in X$, but that $p$ need not be a point of $E$ in
the above definition. Moreover, even if $p \in E$, we may very well have
$f(p) \ne \lim_{x\to p} f(x)$.

\begin{center}
\begin{tikzpicture}[x=1mm, y=1mm]
  % X axis with deleted neighbourhood of p
  \draw[axis] (-2,0) -- (52,0);
  \node[lbl, anchor=north] at (50,-2) {$E \subset X$};
  \node[ptgap] at (24,0) {};
  \node[mlbl, anchor=north] at (24,-2.2) {$p$};
  \draw[seg] (15,0) -- (22.8,0);
  \draw[seg] (25.2,0) -- (33,0);
  \node[ptopen, minimum size=2.4pt] at (15,0) {};
  \node[ptopen, minimum size=2.4pt] at (33,0) {};
  \draw[guide] (15,2) -- (15,7); \draw[guide] (33,2) -- (33,7);
  \node[lbl, anchor=south] at (24,7.5) {punctured $\delta$-ball: $p$ itself excluded};
  % arrow to Y
  \draw[vec] (56,3) -- (66,3);
  \node[mlbl, anchor=south] at (61,4.5) {$f$};
  % Y axis with eps ball around q
  \draw[axis] (70,-6) -- (70,14);
  \node[lbl, anchor=west] at (72,13) {$Y$};
  \node[ptfill, minimum size=2.6pt] at (70,4) {};
  \node[mlbl, anchor=west] at (72,4) {$q$};
  \draw[seg] (70,0.5) -- (70,7.5);
  \node[mlbl, anchor=west] at (72,8.5) {$\varepsilon$-ball};
\end{tikzpicture}
\end{center}
\figcap{The contract of 4.1: however small an $\varepsilon$-ball is drawn
about $q$, some punctured $\delta$-ball about $p$ is mapped entirely into it.
The value $f(p)$ --- if it exists at all --- plays no part.}

We can recast this definition in terms of limits of sequences:

\ritem{4.2}{Theorem}
\emph{Let $X$, $Y$, $E$, $f$, and $p$ be as in Definition 4.1. Then}
\begin{equation}
  \lim_{x\to p} f(x) = q \tag{4.4}
\end{equation}
\emph{if and only if}
\begin{equation}
  \lim_{n\to\infty} f(p_n) = q \tag{4.5}
\end{equation}
\emph{for every sequence $(p_n)$ in $E$ such that}
\begin{equation}
  p_n \ne p, \qquad \lim_{n\to\infty} p_n = p. \tag{4.6}
\end{equation}

\begin{proof}
Suppose (4.4) holds. Choose $(p_n)$ in $E$ satisfying (4.6). Let
$\varepsilon > 0$ be given. Then there exists $\delta > 0$ such that
$d_Y(f(x),q) < \varepsilon$ if $x \in E$ and $0 < d_X(x,p) < \delta$. Also,
there exists $N$ such that $n > N$ implies $0 < d_X(p_n,p) < \delta$. Thus,
for $n > N$, we have $d_Y(f(p_n),q) < \varepsilon$, which shows that (4.5)
holds.

Conversely, suppose (4.4) is false. Then there exists some $\varepsilon > 0$
such that for every $\delta > 0$ there exists a point $x \in E$ (depending on
$\delta$), for which $d_Y(f(x),q) \ge \varepsilon$ but
$0 < d_X(x,p) < \delta$. Taking $\delta_n = 1/n$ $(n = 1,2,3,\dots)$, we thus
find a sequence in $E$ satisfying (4.6) for which (4.5) is
false.\kn{The forward direction chains two contracts: $\varepsilon$ buys a
$\delta$ from $f$, and that $\delta$ buys an $N$ from the sequence. The
converse is proved \emph{contrapositively}, and note its shape: negating
``for every $\varepsilon$ there is $\delta$'' yields one bad $\varepsilon$
whose every $\delta$ fails --- and harvesting one witness $x_n$ for each
$\delta_n = 1/n$ assembles the bad sequence. (Silently, a choice is made for
each $n$; compare the note at 2.12.)}
\end{proof}

\ritem{}{Corollary}
\emph{If $f$ has a limit at $p$, this limit is unique.}

This follows from Theorems 3.2(b) and 4.2.

\begin{deeper}[The bridge, and which way to cross it]
Theorem 4.2 is a dictionary: function limits on one side, sequence limits on
the other. Its value is strategic, and each direction has a distinct use.

\smallskip
\textbf{To import theorems, cross left to right.} All of Chapter 3's limit
arithmetic transfers instantly --- that is the entire proof of 4.4, one line
long. Anything proved for sequences is proved for function limits, free.

\smallskip
\textbf{To refute limits, cross right to left.} To show $\lim_{x\to p} f(x)$
does not exist, produce \emph{two} sequences approaching $p$ along which $f$
tends to different values --- the function-limit analogue of the
two-subsequence trick of 3.1(d). Example 4.27(d) will use exactly this on
$\sin(1/x)$-type behaviour.

\smallskip
The quantifier ``\emph{every} sequence'' is what makes the dictionary
faithful, and it is why the refutation direction is cheap while the
verification direction through sequences is usually hopeless --- one cannot
check all sequences one by one. For verifying, $\varepsilon$--$\delta$ or 4.8
are the tools.
\end{deeper}

\ritem{4.3}{Definition}
Suppose we have two complex functions, $f$ and $g$, both defined on $E$. By
$f+g$ we mean the function which assigns to each point $x$ of $E$ the number
$f(x)+g(x)$. Similarly we define the difference $f-g$, the product $fg$, and
the quotient $f/g$ of two functions, with the understanding that the quotient
is defined only at those points $x$ of $E$ at which $g(x) \ne 0$. If $f$
assigns to each point $x$ of $E$ the same number $c$, then $f$ is said to be a
\emph{constant} function, or simply a \emph{constant}, and we write $f = c$.
If $f$ and $g$ are real functions, and if $f(x) \ge g(x)$ for every
$x \in E$, we shall sometimes write $f \ge g$, for brevity.

Similarly, if $\mathbf{f}$ and $\mathbf{g}$ map $E$ into $\R^k$, we define
$\mathbf{f}+\mathbf{g}$ and $\mathbf{f}\cdot\mathbf{g}$ by
\[ (\mathbf{f}+\mathbf{g})(x) = \mathbf{f}(x)+\mathbf{g}(x), \qquad
   (\mathbf{f}\cdot\mathbf{g})(x) = \mathbf{f}(x)\cdot\mathbf{g}(x); \]
and if $\lambda$ is a real number, $(\lambda\mathbf{f})(x) =
\lambda\mathbf{f}(x)$.

\ritem{4.4}{Theorem}
\emph{Suppose $E \subset X$, a metric space, $p$ is a limit point of $E$, $f$
and $g$ are complex functions on $E$, and}
\[ \lim_{x\to p} f(x) = A, \qquad \lim_{x\to p} g(x) = B. \]
\emph{Then}
\begin{enumerate}[label=(\alph*), leftmargin=2.2em, itemsep=2pt, topsep=3pt]
  \item $\displaystyle\lim_{x\to p} (f+g)(x) = A+B$;
  \item $\displaystyle\lim_{x\to p} (fg)(x) = AB$;
  \item $\displaystyle\lim_{x\to p} \left(\frac fg\right)(x) = \frac AB$,
        \emph{if $B \ne 0$.}
\end{enumerate}

\begin{proof}
In view of Theorem 4.2, these assertions follow immediately from the analogous
properties of sequences (Theorem 3.3).\kn{The dictionary in action: not one
$\varepsilon$ appears, because every $\varepsilon$ was already spent in 3.3.
This is the chapter's characteristic proof --- one sentence of translation
--- and the same sentence will prove 4.9. When a proof is this short, the
work was done elsewhere; the skill is knowing where.}
\end{proof}

\emph{Remark:} If $\mathbf{f}$ and $\mathbf{g}$ map $E$ into $\R^k$, then (a)
remains true, and (b) becomes

\emph{(b$'$)} $\displaystyle\lim_{x\to p}
(\mathbf{f}\cdot\mathbf{g})(x) = \mathbf{A}\cdot\mathbf{B}$.

(Compare Theorem 3.4.)
